% ============================================================
% VERSIÓN INDIVIDUAL — DANIEL DUARTE CORDERO
% ============================================================
\section{Consideraciones Éticas Preliminares}
\label{sec:etica}

\subsection{Sensibilidad de los datos}
\label{subsec:sensibilidad}
Los replays analizados describen partidas entre agentes artificiales y
estados del juego; en la auditoría realizada no se identificaron datos
personales sensibles de jugadores humanos. Sí existen nombres o IDs de
agentes/equipos que se tratarán como metadata de auditoría y no como
features. \textbf{\% TODO integración:} verificar y citar la licencia
exacta con la que The Pokémon Company/Kaggle distribuye los replays y
resumir explícitamente qué usos y redistribución permite. No se infiere
la licencia a partir de la mera disponibilidad pública.

\subsection{Riesgo de sesgo}
\label{subsec:sesgo}
La muestra no es aleatoria respecto del universo de Pokémon TCG. Los
replays de la competencia se seleccionan por rating medio y se truncan
a un volumen diario, por lo que representan principalmente agentes
artificiales competitivos. La auditoría además muestra concentración de
mazos y agentes, así como cambio fuerte entre días: el top-10 de mazos
concentra 39.94\% de las instancias auditadas y las distribuciones de
mazos/agentes cambian sustancialmente entre el primer y último día.
Esto puede favorecer atajos asociados al metajuego dominante y degradar
el desempeño cuando cambian composiciones o versiones.

La literatura revisada refuerza que este riesgo no es meramente
hipotético: VGC-Bench encuentra diferencias entre rendimiento en equipos
vistos y generalización a equipos no vistos, mientras Chitayat et al.
separan evaluación en parches conocidos y posteriores
\cite{angliss_vgcbench_2026,chitayat_meta_2023}.

\subsection{Límites de generalización declarados}
\label{subsec:limitesgeneralizacion}
Los resultados no deben extrapolarse a jugadores humanos, al Pokémon
TCG físico, a todas las estrategias posibles, ni a versiones del motor,
mazos o poblaciones que no hayan sido cubiertos por los protocolos de
evaluación. Tampoco un buen resultado en el split agrupado implica
robustez temporal: esa afirmación solo podrá hacerse si el modelo
mantiene desempeño bajo el protocolo cronológico. De forma análoga, un
buen AUC global no implica que la probabilidad esté calibrada ni que el
modelo funcione igual en etapas tempranas y tardías.

\subsection{Mitigación prevista antes de entrenar}
\label{subsec:mitigacionetica}
\begin{itemize}
  \item Separar y reportar los protocolos agrupado, temporal y por
  mazo/matchup, sin ocultar caídas de desempeño bajo cambio de
  distribución.
  \item Prohibir identidades de agente, ratings y metadata similares
  como features; usar la identidad de mazo solo en una ablación
  controlada y evaluarla con holdout por mazo.
  \item Reportar explícitamente composición de la muestra, versiones,
  concentración y límites de población para evitar que el modelo se
  presente como generalizable a humanos o a todo Pokémon TCG.
\end{itemize}

\subsection{Diferido a la Etapa 3}
\label{subsec:diferido}
La ética del despliegue, los escenarios de uso indebido y los riesgos
operativos del agente inteligente se abordan en la etapa final.
